% GENERATED FILE. Edit canonical lessons, then run npm run generate:books.
% canonical-source-hash: fnv1a64:dd700671895f4f64
% canonical-lessons: KA-C316-maneginta, KA-C316-tingalinda, KA-R316-comparing-and-time

\chapter{Bigger, Best and How Long}
\label{ch:ka-bigger-best-and-how-long}
\hlchaptermodality{316}

\begin{chapteropening}
\textbf{By the end of this chapter:} \emph{I can compare two things with -ginta, say which is most with ellakkinta, say how long something has gone on with -inda, and give a span from … until.}

Complete the last lesson of chapter 316: compare with -ginta and ellakkinta, and say how long with -inda and varege.
\end{chapteropening}

\section[\texorpdfstring{\kn{ಮನೆಗಿಂತ} (maneginta)}{ಮನೆಗಿಂತ (maneginta)}]{\kn{ಮನೆಗಿಂತ} (maneginta) — bigger than, best of all}

\label{lesson:KA-C316-maneginta}



\begin{quote}
Say \textquotedblleft{}big\textquotedblright{}, \textquotedblleft{}small\textquotedblright{} and \textquotedblleft{}more\textquotedblright{}. (\emph{doḍḍa}, \emph{cikka}, \emph{heccu}.) Then \textquotedblleft{}I would love to go to Mysuru one day.\textquotedblright{} (\emph{nanage ondu dina maisūrige hōgabēku anta āse ide}.)
\end{quote}

\subsection*{You'll want to know: -\kn{ಗಿಂತ}, than}
\begin{quote}
\textbf{\kn{ಈ} \kn{ಮನೆ} \kn{ಆ} \kn{ಮನೆಗಿಂತ} \kn{ದೊಡ್ಡದು}.}
\end{quote}

\begin{quote}
\emph{ī mane ā maneginta doḍḍadu.} — \textquotedblleft{}This house is bigger than that house.\textquotedblright{}
\end{quote}

The thing you measure against takes \textbf{-\kn{ಗಿಂತ}}, \textquotedblleft{}than\textquotedblright{}. It is the \textquotedblleft{}to\textquotedblright{} ending with \textbf{-\kn{ಇಂತ}} on it: \textbf{\kn{ಮನೆಗೆ}} $\to$ \textbf{\kn{ಮನೆಗಿಂತ}}; \textbf{\kn{ನನಗೆ}} $\to$ \textbf{\kn{ನನಗಿಂತ}} (\emph{nanaginta}, \textquotedblleft{}than me\textquotedblright{}). Where \textquotedblleft{}to\textquotedblright{} is \textbf{-\kn{ಕ್ಕೆ}}, \textquotedblleft{}than\textquotedblright{} is \textbf{-\kn{ಕ್ಕಿಂತ}}: \textbf{\kn{ಮರಕ್ಕೆ}} $\to$ \textbf{\kn{ಮರಕ್ಕಿಂತ}} (\emph{marakkinta}, \textquotedblleft{}than the tree\textquotedblright{}).

The adjective needs no \textquotedblleft{}-er\textquotedblright{}. \textbf{\kn{ದೊಡ್ಡ}} stays \textbf{\kn{ದೊಡ್ಡ}}; at the end of the sentence it takes \textbf{-\kn{ದು}}: \textbf{\kn{ದೊಡ್ಡದು}} (\emph{doḍḍadu}, \textquotedblleft{}(it is) big\textquotedblright{}), \textbf{\kn{ಚಿಕ್ಕದು}} (\emph{cikkadu}, \textquotedblleft{}(it is) small\textquotedblright{}), like the \textbf{\kn{ಒಳ್ಳೆಯದು}} of \textquotedblleft{}it is good\textquotedblright{}.

\begin{quote}
\textbf{\kn{ಬಸ್ಸಿಗಿಂತ} \kn{ರೈಲು} \kn{ಬೇಗ} \kn{ಹೋಗುತ್ತದೆ}.}
\end{quote}

\begin{quote}
\emph{bassiginta railu bēga hōguttade.} — \textquotedblleft{}The train goes faster than the bus.\textquotedblright{}
\end{quote}

\begin{quote}
\textbf{\kn{ಕಾಫಿಗಿಂತ} \kn{ಚಹಾ} \kn{ನನಗೆ} \kn{ಹೆಚ್ಚು} \kn{ಇಷ್ಟ}.}
\end{quote}

\begin{quote}
\emph{kāphiginta cahā nanage heccu iṣṭa.} — \textquotedblleft{}I like tea more than coffee.\textquotedblright{}
\end{quote}

With a liking, \textbf{\kn{ಹೆಚ್ಚು}}, \textquotedblleft{}more\textquotedblright{}, says how much more.

\subsection*{You'll want to know: \kn{ಎಲ್ಲಕ್ಕಿಂತ}, most of all}
\begin{quote}
\textbf{\kn{ಇದು} \kn{ಎಲ್ಲಕ್ಕಿಂತ} \kn{ಚಿಕ್ಕ} \kn{ಕೋಣೆ}.}
\end{quote}

\begin{quote}
\emph{idu ellakkinta cikka kōṇe.} — \textquotedblleft{}This is the smallest room.\textquotedblright{}
\end{quote}

\textbf{\kn{ಎಲ್ಲಕ್ಕಿಂತ}} is \textbf{\kn{ಎಲ್ಲ}}, \textquotedblleft{}all\textquotedblright{}, the \textbf{\kn{ಎಲ್ಲಾ}} you know, with \textbf{-\kn{ಕ್ಕಿಂತ}}: \textquotedblleft{}than all\textquotedblright{}, so \textquotedblleft{}most\textquotedblright{}. For people it is \textbf{\kn{ಎಲ್ಲರಿಗಿಂತ}} (\emph{ellariginta}, \textquotedblleft{}than everyone\textquotedblright{}):

\begin{quote}
\textbf{\kn{ಅವಳು} \kn{ಎಲ್ಲರಿಗಿಂತ} \kn{ಬೇಗ} \kn{ಬಂದಳು}.}
\end{quote}

\begin{quote}
\emph{avaḷu ellariginta bēga bandaḷu.} — \textquotedblleft{}She came earliest of all.\textquotedblright{}
\end{quote}

\subsection*{Your turn}
Say these aloud:

\begin{itemize}
\raggedright
  \item \textquotedblleft{}This house is bigger than that house\textquotedblright{} — \emph{ī mane ā maneginta doḍḍadu}
  \item \textquotedblleft{}The train goes faster than the bus\textquotedblright{} — \emph{bassiginta railu bēga hōguttade}
  \item \textquotedblleft{}I like tea more than coffee\textquotedblright{} — \emph{kāphiginta cahā nanage heccu iṣṭa}
  \item \textquotedblleft{}This is the smallest room\textquotedblright{} — \emph{idu ellakkinta cikka kōṇe}
\end{itemize}

\begin{tcolorbox}[breakable,colback=teal!4,colframe=teal!35!black,title={Before you move on}]
\textquotedblleft{}Than the house\textquotedblright{}? (\textbf{\kn{ಮನೆಗಿಂತ}}.) \textquotedblleft{}Than me\textquotedblright{}? (\textbf{\kn{ನನಗಿಂತ}}.) \textquotedblleft{}It is bigger\textquotedblright{}? (\textbf{\emph{}\kn{ದೊಡ್ಡದು}\emph{\textbf{\emph{: no \textquotedblleft{}-er\textquotedblright{}.) \textquotedblleft{}Most of all\textquotedblright{}? (\textbf{\kn{ಎಲ್ಲಕ್ಕಿಂತ}}; for people, \textbf{\kn{ಎಲ್ಲರಿಗಿಂತ}}.)}}}}
\end{tcolorbox}
\section[\texorpdfstring{\kn{ಎರಡು} \kn{ತಿಂಗಳಿಂದ} (eraḍu tiṅgaḷinda)}{ಎರಡು ತಿಂಗಳಿಂದ (eraḍu tiṅgaḷinda)}]{\kn{ಎರಡು} \kn{ತಿಂಗಳಿಂದ} (eraḍu tiṅgaḷinda) — how long}

\label{lesson:KA-C316-tingalinda}



\begin{quote}
Say \textquotedblleft{}time\textquotedblright{} (\emph{samaya}) and \textquotedblleft{}until\textquotedblright{} (\emph{varege}). Then \textquotedblleft{}This house is bigger than that house.\textquotedblright{} (\emph{ī mane ā maneginta doḍḍadu}.)
\end{quote}

\subsection*{You'll want to know: -\kn{ಇಂದ}, for … now}
\begin{quote}
\textbf{\kn{ನಾನು} \kn{ಎರಡು} \kn{ತಿಂಗಳಿಂದ} \kn{ಕನ್ನಡ} \kn{ಕಲಿಯುತ್ತಿದ್ದೇನೆ}.}
\end{quote}

\begin{quote}
\emph{nānu eraḍu tiṅgaḷinda kannaḍa kaliyuttiddēne.} — \textquotedblleft{}I have been learning Kannada for two months.\textquotedblright{}
\end{quote}

\textbf{\kn{ತಿಂಗಳು}} (\emph{tiṅgaḷu}, \textquotedblleft{}a month\textquotedblright{}) with \textbf{-\kn{ಇಂದ}}, \textquotedblleft{}from\textquotedblright{}: \textquotedblleft{}from two months (back to now)\textquotedblright{}. The \textbf{-\kn{ಇಂದ}} of \textbf{\kn{ಇಲ್ಲಿಂದ}}, \textquotedblleft{}from here\textquotedblright{}, now measures time. The verb stays in the present, because the learning is still going on: \textbf{\kn{ಕಲಿಯುತ್ತಿದ್ದೇನೆ}} is \textbf{\kn{ಕಲಿಯುತ್ತ}} + \textbf{\kn{ಇದ್ದೇನೆ}}, \textquotedblleft{}I am, learning\textquotedblright{}, the present twin of \textbf{\kn{ಓದುತ್ತಿದ್ದೆ}}. English says \textquotedblleft{}have been\textquotedblright{}; Kannada says \textquotedblleft{}am\textquotedblright{}.

\begin{quote}
\textbf{\kn{ನೀವು} \kn{ಎಷ್ಟು} \kn{ಸಮಯದಿಂದ} \kn{ಇಲ್ಲಿ} \kn{ಇದ್ದೀರಿ}?}
\end{quote}

\begin{quote}
\emph{nīvu eṣṭu samayadinda illi iddīri?} — \textquotedblleft{}How long have you been here?\textquotedblright{}
\end{quote}

\textbf{\kn{ಇದ್ದೀರಿ}} (\emph{iddīri}, \textquotedblleft{}you are\textquotedblright{}), said with respect. The answer: \textbf{\kn{ಮೂರು} \kn{ವರ್ಷದಿಂದ}} (\emph{mūru varṣadinda}, \textquotedblleft{}for three years\textquotedblright{}).

A stretch that is over takes no ending at all: \textbf{\kn{ನಾನು} \kn{ಒಂದು} \kn{ವಾರ} \kn{ಮೈಸೂರಿನಲ್ಲಿ} \kn{ಇದ್ದೆ}} (\emph{nānu ondu vāra maisūrinalli idde}, \textquotedblleft{}I was in Mysuru for a week\textquotedblright{}). \textbf{\kn{ವಾರ}} (\emph{vāra}, \textquotedblleft{}a week\textquotedblright{}).

\subsection*{You'll want to know: from … until}
\begin{quote}
\textbf{\kn{ನಾನು} \kn{ಬೆಳಗ್ಗೆ} \kn{ಒಂಬತ್ತು} \kn{ಗಂಟೆಯಿಂದ} \kn{ಸಂಜೆ} \kn{ಐದು} \kn{ಗಂಟೆಯವರೆಗೆ} \kn{ಕೆಲಸ} \kn{ಮಾಡುತ್ತೇನೆ}.}
\end{quote}

\begin{quote}
\emph{nānu beḷagge ombattu gaṇṭeyinda sañje aidu gaṇṭeyavarege kelasa māḍuttēne.} — \textquotedblleft{}I work from nine in the morning until five in the evening.\textquotedblright{}
\end{quote}

\textbf{-\kn{ಇಂದ}} marks where the stretch starts, \textbf{\kn{ವರೆಗೆ}} where it ends: \textbf{\kn{ಗಂಟೆಯಿಂದ}} … \textbf{\kn{ಗಂಟೆಯವರೆಗೆ}}. The same pair works for places: \textbf{\kn{ಇಲ್ಲಿಂದ} \kn{ನಿಲ್ದಾಣದವರೆಗೆ}} (\emph{illinda nildāṇadavarege}, \textquotedblleft{}from here to the station\textquotedblright{}).

\subsection*{Your turn}
Say these aloud:

\begin{itemize}
\raggedright
  \item \textquotedblleft{}I have been learning Kannada for two months\textquotedblright{} — \emph{nānu eraḍu tiṅgaḷinda kannaḍa kaliyuttiddēne}
  \item \textquotedblleft{}How long have you been here?\textquotedblright{} — \emph{nīvu eṣṭu samayadinda illi iddīri?}
  \item \textquotedblleft{}For three years\textquotedblright{} — \emph{mūru varṣadinda}
  \item \textquotedblleft{}I work from nine in the morning until five in the evening\textquotedblright{} — \emph{nānu beḷagge ombattu gaṇṭeyinda sañje aidu gaṇṭeyavarege kelasa māḍuttēne}
\end{itemize}

\begin{tcolorbox}[breakable,colback=teal!4,colframe=teal!35!black,title={Before you move on}]
\textquotedblleft{}For two months now\textquotedblright{}? (\textbf{\kn{ಎರಡು} \kn{ತಿಂಗಳಿಂದ}}.) Which tense goes with it? (\textbf{The present}, for what is still going on.) \textquotedblleft{}How long?\textquotedblright{} (\textbf{\kn{ಎಷ್ಟು} \kn{ಸಮಯದಿಂದ}?}) \textquotedblleft{}From … until\textquotedblright{}? (\textbf{-\kn{ಇಂದ}} … \textbf{\kn{ವರೆಗೆ}}.)
\end{tcolorbox}
\section[(review)]{Review — comparing, and how long}

\label{lesson:KA-R316-comparing-and-time}



\begin{quote}
Read \textbf{\kn{ಸ್ವಚ್ಛ}} and \textbf{\kn{ಔಷಧ}} aloud. (\emph{svaccha}, \emph{auṣadha}.) Name the letter each one taught you. (\textbf{\kn{ಛ}}; \textbf{\kn{ಔ}}.)
\end{quote}

\subsection*{Your turn}
Say these aloud:

\begin{itemize}
\raggedright
  \item My town is smaller than Mysuru.
  \item I like the train more than the bus. One ticket to Mysuru, please.
  \item This is the biggest shop in our town.
  \item I have been learning Kannada for two months, and I am glad.
  \item I was in Mysuru from Monday until Friday.
\end{itemize}

\begin{tcolorbox}[breakable,colback=teal!4,colframe=teal!35!black,title={Before you move on}]
\textquotedblleft{}Than the bus\textquotedblright{}? (\textbf{\kn{ಬಸ್ಸಿಗಿಂತ}}.) \textquotedblleft{}Most of all\textquotedblright{}? (\textbf{\kn{ಎಲ್ಲಕ್ಕಿಂತ}}.) \textquotedblleft{}For three years now\textquotedblright{}? (\textbf{\kn{ಮೂರು} \kn{ವರ್ಷದಿಂದ}}.) \textquotedblleft{}From nine until five\textquotedblright{}? (\textbf{\kn{ಒಂಬತ್ತು} \kn{ಗಂಟೆಯಿಂದ} \kn{ಐದು} \kn{ಗಂಟೆಯವರೆಗೆ}}.)
\end{tcolorbox}
